\documentclass{tufte-handout}

\usepackage{fancyvrb}
\fvset{fontsize=\normalsize}

\usepackage{hyperref}
\hypersetup{colorlinks=true}

\usepackage{graphicx}

% Hide page numbers in ToC
\let\Contentsline\contentsline
\renewcommand\contentsline[3]{\Contentsline{#1}{#2}{}}

\title{OpenStack Orchestration with Heat}
\author{\href{mailto:zbitter@redhat.com}{Zane Bitter}}
\date{24 October 2013}

\usepackage{pdfpages}

\begin{document}
\maketitle

\marginnote{\tableofcontents \vspace{2em}}

\begin{abstract}
An introduction to the OpenStack Orchestration project, Heat, and an explanation of how orchestration can help simplify the deployment and management of your cloud application by allowing you to represent infrastructure as code. Autoscaling, another service provided by Heat, is also covered.
\end{abstract}

\section{What is Heat?}

Heat is the project providing Orchestration for OpenStack. It allows cloud infrastructure to be represented in a declarative form (`template'). The Heat service parses these templates and calls the various OpenStack APIs to create or modify the deployment, while maintaining correct dependencies between them. The collection of resources associated with a template is known as a stack. One of the resource types available is a stack, so templates can be composed into hierarchical structures.

\begin{figure}[b!]
\centering
\includegraphics[trim=3em 6em 3em 10em,width=\textwidth]{openstack-arch.pdf}
\caption{OpenStack architecture from a user's perspective, showing how Orchestration interacts with other services.}
\label{fig:openstack-arch}
\end{figure}

Heat provides equivalent functionality to the AWS CloudFormation service. CloudFormation templates are represented as \textsc{Json} documents referencing resources, and Heat is compatible with the CloudFormation template syntax (and a \textsc{Yaml} variant for easier consumption by humans).  Heat also provides a number of Cloudformation-compatible resource implementations, and a CloudFormation-compatible API.

Heat also provides a growing number of OpenStack native resource types, which aim to allow flexible and transparent configuration of the underlying OpenStack resources. Heat also has its own OpenStack-native \textsc{Rest} API, and work is in progress adding functionality to a native template language (the Heat Orchestration Template format, aka ``\textsc{Hot}'').


\section{Heat Templates}

\begin{marginfigure}
\begin{verbatim}
Parameters:
  ssh_key_name:
    Type: String
    Description: ssh keypair name
  image_name:
    Type: String
    Description: The image to boot

Resources:
  my_server:
    Type: OS::Nova::Server
    Properties:
      flavor: m1.small
      key_name: {"Ref": "ssh_key"}
      block_device_mapping:
        device_name: vda
        volume_id: {"Ref": "my_vol"}

  my_vol:
    Type: OS::Cinder::Volume
    Properties:
      size: 20
      image: {"Ref": "image_name"}

Outputs:
  server_ip:
    Description: The server IP
    Value: {"Fn::GetAtt":
                ["my_server",
                 "first_address"]}
\end{verbatim}

\caption{A simple example template that creates a Nova server and an attached Cinder volume. The user can specify the image to boot from and the name of the ssh public key to install on the server, and the IP address of the server is available as an output.}
\label{fig:example-template}
\end{marginfigure}

Heat templates usually take the form of simple \textsc{Yaml} documents. We chose this over \textsc{Json}, which CloudFormation uses, because it is much easier for humans to read and write---and a diff between two versions of a \textsc{Yaml} template is usually trivial to interpret. \textsc{Yaml} is a strict superset of \textsc{Json} though, so \textsc{Json} is still fully supported and templates can be converted between the two formats with no loss of fidelity.

Other than the serialisation format, Heat (for now) hews closely to the CloudFormation template model. (An example template is shown in Figure~\ref{fig:example-template}.) A template has four key elements:

\begin{itemize}
\item An optional \textsc{Parameters} section, which allows user-definable inputs to be specified.
\item An optional \textsc{Mappings} section, which allows key/value lookup of predefined constants.
\item A mandatory \textsc{Resources} section, which describes the resources and relationships which define the application, configuration and infrastructure.
\item An optional \textsc{Outputs} section, which describes the output values to be returned to the user.
\end{itemize}

All resources have a common interface:

\begin{itemize}
\item A number of optional or mandatory \textsc{Properties} which specify inputs that affect how the resource is configured.
\item A number of output attributes, which may be referenced elsewhere in the template using the \texttt{Fn::GetAtt} function.
\end{itemize}

One of the resource types available is a Heat stack, so multiple templates can be composed into hierarchical structures.

A resource may reference any other resource or its attributes, which causes Heat to infer an ordering dependency between them. The template may also specify explicit dependencies where necessary. Lifecycle operations\footnote{That is to say, operations such as creating, deleting or updating a stack.} occur in parallel to the maximum extent possible given the dependencies between resources.

A number of \textsc{Functions} are available which allow some limited manipulation of data within the template, for example to combine parameter data with static data in the template.


\section{Providers and Environments}

The OpenStack ecosystem is obviously far more diverse than the Amazon walled garden, and for that reason Heat has added features in the Havana release cycle aimed at making templates more portable across clouds.

The first issue to deal with is that users are at the mercy of their cloud providers for native (Python) resource type plugins. The Providers feature helps to mitigate that, by allowing users to effectively define their own resource types. Provider resources are simply constructed out of native resources and are expressed using the exact same template syntax as stacks. (In fact, they are implemented as stacks.)

Heat has also added environments, to make providers easier to consume. An environment allows the user to override parts of the template---for example, by specifying a provider template to implement a particular resource type. In the ideal scenario, an organisation could keep both a library of templates (one per application) and an orthogonal library of environments (one per cloud). By selecting the appropriate template and environment, any application could be launched on any cloud without modification.

For example, you may be deploying your application to an OpenStack cloud that features a DBaaS (and the appropriate Heat plugin for it), but testing on a local OpenStack installation with just the basics. You could define an environment for your test cloud that specifies that the DBaaS resource is supplied by a provider template containing a Nova server running MySQL. That allows you to test your application on your local cloud without DBaaS and still deploy the self same template to staging or production.

Providers are a powerful abstraction, and when they are available in the wild it is likely we will see more adventurous uses of them. You can supply a provider template to define a resource type even when a native plugin for that resource type exists. That allows you, for example, to create a provider template for a server resource that is specialised for a particular configuration management system or PaaS. Thus, in the main template you would only need to supply configuration data for your configuration management tool rather than configure the tool itself on every server.


\section{Matahari Roadmap}

A detailed list of what tasks are currently being worked on can be found on the project \href{https://github.com/matahari/matahari/wiki/Backlog}{backlog}. 

\subsection{More Information}

Here is a list of other places you can go to get more information and participate in the project:

\begin{itemize}
\item The project home page is \url{http://www.matahariproject.org/}.
\item You can find the code on \href{https://github.com/matahari/matahari}{github}.
\item You can follow patches and project related discussion on the \href{https://fedorahosted.org/mailman/listinfo/matahari}{mailing list}. 
\end{itemize}


\end{document}
